\chapter{SYSTEM ANALYSIS }

\section{Module Description}
The "Alumni System" is a comprehensive web networking platform designed specifically to connect students, alumni, departments, and administrators. The system is divided into multiple functional modules to ensure a seamless experience:
\begin{itemize}
    \item \textbf{Authentication \& Authorization Module:} Handles user registration, secure login, password resets, and role-based access control (Admin, Alumni, Student, Department). It ensures only verified users gain access to the platform's core features.
    \item \textbf{Profile Management Module:} Allows users to create detailed profiles containing their personal information, academic background (department, graduation year, roll number), skills, current employment, and social links (like LinkedIn). It supports profile pictures and varied layouts dynamically adjusting based on the user's role.
    \item \textbf{Networking \& Feed Module:} The core engaging feature where users can create posts (including text, images, and videos) visible to their network. Other users can interact with these posts by "Liking" them or leaving comments, similar to professional networks like LinkedIn.
    \item \textbf{Direct Messaging Module:} Provides a real-time private chat interface allowing direct communication between any two users in the system. It tracks read/unread statuses and displays conversations in an intuitive interface.
    \item \textbf{Notification Module:} A sophisticated universal notification engine that polls system activities and alerts users about relevant updates. It triggers database notifications for new direct messages and new network posts, accessible instantly via a top-navigation dropdown menu.
    \item \textbf{Search \& Discovery Module:} Allows users to find and connect with peers by searching for specific alumni or filtering the network based on departments, graduation years, or skills.
    \item \textbf{Department \& Admin Module:} Provides administrative oversight. Admins can approve or reject new user registrations to maintain an exclusive network, while Departments have their own dedicated interfaces to view alumni belonging to their specific branches.
\end{itemize}

\subsection{Module Summary Table}
\begin{table}[H]
    \centering
    \begin{tabular}{|l|p{10cm}|}
        \hline
        \textbf{Module Name} & \textbf{Primary Purpose} \\ \hline
        Authentication & Secure user access and role validation. \\ \hline
        Profile Management & User biography, academic, and skills data. \\ \hline
        Networking Feed & Public updates, posts, and interactive comments. \\ \hline
        Direct Messaging & Real-time private 1-on-1 communication. \\ \hline
        Notifications & Universal alerts for new messages and posts. \\ \hline
        Admin/Department & User moderation and alumni tracking interfaces. \\ \hline
    \end{tabular}
    \caption{Summary of System Modules}
    \label{tab:modules}
\end{table}

\section{Business Rules}
Business rules are specific guidelines that govern the operations of the Alumni System to ensure consistency, security, and alignment with institutional goals.\\\\
\textbf{User Registration \& Verification}
\begin{itemize}
    \item All new user registrations (except initial Super Admins) require administrative approval before they are fully active in the system.
    \item A user must complete their basic Profile Setup—including required fields like Department and Roll Number—before accessing the main dashboard and networking features.
\end{itemize}

\textbf{Content \& Interactions}
\begin{itemize}
    \item Only verified users can create posts, like existing posts, or leave comments.
    \item Users can only delete posts that they have authored themselves. Administrators hold global privileges to delete any inappropriate posts across the network.
    \item Multimedia uploads are strictly validated: images must be formatted properly and not exceed 10MB, while videos must be under 50MB.
\end{itemize}

\textbf{Networking \& Privacy}
\begin{itemize}
    \item Private messages can only be sent between registered and active users.
    \item A read-receipt system ensures users know when their messages have been seen.
    \item Certain profile fields, such as Roll Numbers and specific Academic details, are deliberately hidden for users holding 'Department' and 'Admin' roles to ensure the interface remains clean and relevant to their status.
\end{itemize}

\newpage
\section{Analysis of Architecture (and/ or) Algorithms}
{\large \textbf{Laravel PHP Framework}}\\\\
The Alumni System is built on Laravel, a robust and elegant PHP web framework. Laravel provides the foundation for secure routing, MVC architecture, and database interactions, ensuring a scalable and maintainable networking platform.\\\\
\begin{enumerate}
    \item \textbf{MVC Architecture:}
    \begin{itemize}
        \item \textbf{Models:} Eloquent ORM is used for seamless database interactions mapping directly to tables like \texttt{users}, \texttt{profiles}, \texttt{posts}, \texttt{messages}, and \texttt{notifications}.
        \item \textbf{Views:} Blade templating engine renders dynamic UI components safely and efficiently.
        \item \textbf{Controllers:} Handle the core business logic, from authenticating users to processing new post submissions and real-time messaging data.
    \end{itemize}
    \item \textbf{Database Migrations \& Relationships:}
    \begin{itemize}
        \item Laravel’s migration system ensures the PostgreSQL database schema is version-controlled.
        \item Complex eloquent relationships (e.g., User \texttt{hasOne} Profile, User \texttt{hasMany} Posts) allow for rapid data retrieval, such as loading a user's entire feed of posts alongside their profile picture.
    \end{itemize}
    \item \textbf{Security:}
    \begin{itemize}
        \item Built-in CSRF protection safeguards all forms.
        \item Passwords are cryptographically hashed, and SQL injection is prevented automatically through Eloquent's prepared statements.
    \end{itemize}
\end{enumerate}

\textbf{Why Laravel for the Alumni System?}\\\\
Laravel provides a comprehensive web development framework that ensures rapid deployment, high security, and elegant code structure. Its built-in ecosystem handles complex requirements like Authentication and Notifications out-of-the-box, making it the ideal choice for a feature-rich social network.

\section{Feasibility Analysis}

\subsection{Technical Feasibility}
Technical feasibility assesses whether a project can be successfully developed and implemented using the current available technology, skills, and resources.\\\\
\noindent
The chosen tech stack for the Alumni System—Laravel (PHP), TailwindCSS, and PostgreSQL—is highly feasible and aligns perfectly with modern web development standards. The backend infrastructure easily supports complex relational data required for social networking (users, posts, likes, messages). The frontend utilizes TailwindCSS for a responsive, modern, "LinkedIn-style" UI that works flawlessly across desktop and mobile devices.

\subsection{Economical Feasibility}
Economic feasibility refers to the assessment of the financial aspects of a project to determine whether it is cost-effective and viable.\\\\
\noindent
The project demonstrates strong economic viability. By utilizing free, open-source technologies (PHP, Laravel, PostgreSQL, TailwindCSS), licensing costs are virtually eliminated. The application can be hosted on very affordable cloud providers (like DigitalOcean or AWS EC2) using simple standardized deployment pipelines. Future maintenance is cost-effective due to Laravel’s vast developer community and extensive documentation.

\subsection{Operational Feasibility}
Operational feasibility evaluates whether the project aligns with operational requirements and whether it can be smoothly integrated into existing processes.\\\\
\noindent
The system is operationally sound. It replaces manual alumni tracking and networking with an automated, user-friendly digital platform. Features like Admin Approval for new accounts ensure the network remains secure and exclusive to actual students and graduates. The dynamic dashboard, instant search, and real-time notifications ensure users remain engaged and can operate the platform intuitively without prior training.

\section{System Environment}

\subsection{Software Environment}
\textbf{Backend Framework:}\\
Laravel 11.x (PHP 8.2+) serves as the core framework for logic, routing, and database interaction.\\\\
\textbf{Frontend Styling:}\\
Tailwind CSS, Blade Templates, minimal Vanilla JavaScript for interactivity (like dropdowns and modals).\\\\
\textbf{Database Management:}\\
PostgreSQL for relational data mapping of users, messages, and posts.\\\\
\textbf{Development Environment:}\\
Artisan Serve and Vite (for asset bundling and hot-module replacement). Git and GitHub are used for version control.

\subsection{Hardware Environment}
\textbf{Processor (CPU):} 2 Cores (e.g., Intel Xeon or AWS T3)\\\\
\textbf{Memory (RAM):} 4 GB minimum\\\\
\textbf{Storage:} 50 GB SSD (to accommodate user uploads like Profile Pictures, Post Images, and Videos)\\\\
\textbf{Operating System:} Ubuntu 22.04 LTS (or similar Linux distribution)

\newpage

\section{Actors and roles}
The system involves four primary actors:\\\\
\noindent
\textbf{Actors name: Admin}\\\\
\textbf{Role:} 
\begin{itemize}
    \item Has full oversight of the platform. Their primary responsibility is approving or rejecting new user registrations and moderating content (deleting inappropriate posts).
\end{itemize}
\noindent
\textbf{Actors name: Alumni}\\\\
\textbf{Role:} 
\begin{itemize}
    \item Graduates of the institution. They can create posts, update their employment profiles, message other users, and search the network for connections.
\end{itemize}
\noindent
\textbf{Actors name: Student}\\\\
\textbf{Role:} 
\begin{itemize}
    \item Current attendees of the institution. They use the platform to seek mentorship, view alumni careers, and stay updated via the global feed.
\end{itemize}
\noindent
\textbf{Actors name: Department}\\\\
\textbf{Role:} 
\begin{itemize}
    \item Represents specific academic faculty branches. They can view directories of all alumni that graduated from their specific department to track placements and success.
\end{itemize}
